\documentclass[11pt]{article}
\usepackage[margin=1in]{geometry}
\usepackage{amsmath}
\usepackage{array}
\usepackage{hyperref}

\title{Interactive Energy-Aware Physarum Slime Mold Simulation}
\author{Slime Mold WebAssembly Simulation}
\date{}

\begin{document}
\maketitle

\begin{abstract}
This document describes a browser-based simulation of \emph{Physarum polycephalum}
for interactive exploration of food search, energy-aware movement, trail-mediated
collective behavior, and adaptive nutrient transport networks. The implementation
uses C++20 for the simulation core, WebAssembly for browser execution, and a
TypeScript canvas interface for rendering and parameter control.
\end{abstract}

\section{Model Assumptions}
The model is not a full biological reconstruction. It combines a Jones-style
particle approximation of slime mold morphology with a simplified Physarum-inspired
transport graph. Agents are point particles with local sensing, energy budgets,
food consumption, death, and reproduction. The trail field approximates extracellular
chemoattractant and tube reinforcement. Food scent is a diffuse attractor field.
Walls are represented by a repellent field. A low-resolution graph is extracted
from high-trail grid cells and food nodes, then updated with approximate pressure,
flow, and conductivity dynamics.

\section{Variables and Parameters}
\begin{center}
\begin{tabular}{>{\raggedright\arraybackslash}p{0.25\linewidth}
                >{\raggedright\arraybackslash}p{0.62\linewidth}}
\hline
Symbol or parameter & Meaning \\
\hline
$p_i$ & Agent $i$ position in the two-dimensional grid. \\
$\theta_i$ & Agent heading angle. \\
$E_i$ & Agent energy. \\
$T(x,y)$ & Slime trail or chemoattractant field. \\
$F(x,y)$ & Food scent field. \\
$R(x,y)$ & Wall or repellent field. \\
$d_s$ & Sensor distance. \\
$\alpha$ & Sensor angular offset. \\
$w_T,w_F,w_R$ & Trail, food, and repellent sensing weights. \\
$v_{\min},v_{\max}$ & Minimum and maximum agent speed. \\
$C_j$ & Remaining calories at food source $j$. \\
$q_i$ & Trail deposited by agent $i$. \\
$D_T,\lambda_T$ & Trail diffusion and decay coefficients. \\
$D_{ij}$ & Graph edge conductivity. \\
$Q_{ij}$ & Nutrient flow along graph edge $i,j$. \\
$P_i$ & Pressure at graph node $i$. \\
$b_i$ & Source or sink term at graph node $i$. \\
\hline
\end{tabular}
\end{center}

\section{Agent Chemotaxis Model}
Each agent has a heading unit vector
\begin{equation}
\label{eq:heading}
u(\theta) = (\cos \theta, \sin \theta).
\end{equation}
Three sensors are placed in front of the agent:
\begin{align}
\label{eq:sensors}
x_F &= p_i + d_s u(\theta_i),\\
x_L &= p_i + d_s u(\theta_i + \alpha),\\
x_R &= p_i + d_s u(\theta_i - \alpha).
\end{align}
For $k \in \{L,F,R\}$, the sampled sensor score is
\begin{equation}
\label{eq:sensor-score}
S_k = w_T T(x_k) + w_F F(x_k) - w_R R(x_k) + \eta_k,
\end{equation}
where $\eta_k$ is a small random perturbation. If the front sensor is strongest,
the agent keeps its heading. If the front sensor is weakest, the agent turns
randomly left or right. Otherwise it turns toward the stronger lateral sensor.
The position update is
\begin{equation}
\label{eq:position-update}
p_i(t + \Delta t) = p_i(t) + v_i(E_i)\Delta t\,u(\theta_i).
\end{equation}
Speed depends on energy:
\begin{equation}
\label{eq:energy-speed}
v_i(E_i) = v_{\min} + (v_{\max} - v_{\min})
\operatorname{clamp}\left(\frac{E_i}{E_{\max}},0,1\right).
\end{equation}

\section{Food Attraction and Consumption}
Food applies a hunger-weighted directional bias. The raw food force is
\begin{equation}
\label{eq:food-force}
F^{\mathrm{food}}_i =
\sum_j c_j \frac{p^{\mathrm{food}}_j - p_i}
{\lVert p^{\mathrm{food}}_j - p_i \rVert^2 + \epsilon},
\end{equation}
where $c_j$ is proportional to remaining calories and attractor strength.
The heading is mixed toward the force angle with
\begin{equation}
\label{eq:food-beta}
\beta_i =
\operatorname{clamp}\left(
\lVert F^{\mathrm{food}}_i \rVert
\left(1 - \frac{E_i}{E_{\max}}\right)^\gamma,0,1
\right).
\end{equation}
If an agent is inside a food radius, it consumes
$\min(C_j,\mathrm{eatRate}\Delta t)$ calories and converts them to energy with
the food efficiency parameter.

\section{Energy Budget}
Agent energy is updated by
\begin{equation}
\label{eq:energy-budget}
E_i(t+\Delta t) =
\operatorname{clamp}\left(
E_i(t)
- c_{\mathrm{base}}\Delta t
- c_{\mathrm{move}}v_i\Delta t
- c_{\mathrm{sensor}}n_s\Delta t
- c_{\mathrm{trail}}q_i\Delta t
+ E^{\mathrm{food}}_i,
0,E_{\max}
\right).
\end{equation}
Here $n_s=3$ is the number of sensors. Trail deposition is energy-scaled:
\begin{equation}
\label{eq:trail-deposit}
q_i = q_{\mathrm{base}}\operatorname{clamp}\left(\frac{E_i}{E_{\max}},0,1\right).
\end{equation}

\section{Growth, Reproduction, Dormancy, and Death}
If $E_i > E_{\mathrm{split}}$, an agent may divide with probability
$p_{\mathrm{split}}\Delta t$. Parent and child split the energy according to
a split ratio, and the child receives a perturbed heading. If energy drops below
$E_{\mathrm{low}}$, speed and trail secretion are reduced to model dormancy. If
energy remains below $E_{\mathrm{dead}}$ longer than the death time, the agent is
marked dead and no longer participates in the simulation.

\section{Search-versus-Exploit Decision}
The probability of continuing exploratory search is
\begin{equation}
\label{eq:p-search}
P^{\mathrm{search}}_i =
\sigma\left(
a_0
+ a_E\frac{E_i}{E_{\max}}
- a_F\,\mathrm{localFood}_i
- a_T\,\mathrm{localTrail}_i
+ a_N\,\mathrm{novelty}_i
- a_C\,\mathrm{localCost}_i
\right),
\end{equation}
where $\sigma(z)=1/(1+e^{-z})$. Search mode increases random exploration and
reduces food/trail exploitation; exploit mode does the reverse.

\section{Trail Deposition, Diffusion, and Decay}
The trail field evolves according to
\begin{equation}
\label{eq:trail-diffusion}
T(t+\Delta t) =
(1-\lambda_T\Delta t)\left(T(t)+D_T\Delta t\nabla^2 T\right).
\end{equation}
The implementation uses an explicit five-point Laplacian with internal substeps
when diffusion and time step combinations would otherwise be unstable. Trail values
are clamped to $[0,T_{\max}]$, and wall cells force trail to zero.

\section{Nutrient Transport Graph}
A low-resolution graph is sampled every \texttt{graphStride} grid cells. Cells with
$T(x,y)$ above the trail threshold become active nodes. Food sources are always
inserted as source nodes, and the live colony center is inserted as a sink node.
Adjacent active grid nodes become graph edges with length $L_{ij}$ and conductivity
$D_{ij}$.

\section{Flow-Conductivity Feedback}
Edge flow is modeled as
\begin{equation}
\label{eq:flow}
Q_{ij} = \frac{D_{ij}}{L_{ij}+\epsilon}(P_i-P_j).
\end{equation}
At each node, approximate conservation is
\begin{equation}
\label{eq:conservation}
\sum_j Q_{ij} = b_i.
\end{equation}
The pressure system is solved with Gauss-Seidel relaxation. Conductivity follows
\begin{equation}
\label{eq:conductivity-ode}
\frac{dD_{ij}}{dt} = \alpha_D |Q_{ij}|^\mu - \lambda_D D_{ij},
\end{equation}
with the discrete update clamped to configured minimum and maximum conductivities.
High-flow paths thicken; low-flow paths decay.

\section{Shortest Nutrient Path Optimization}
For shortest nutrient paths, Dijkstra uses the edge weight
\begin{equation}
\label{eq:path-weight}
w_{ij} = \frac{L_{ij}}{D_{ij}+\epsilon}
+ \lambda_{\mathrm{energy}}\mathrm{movementCost}_{ij}.
\end{equation}
The displayed highlighted path is the lowest-cost path from food nodes to the
colony sink. Transport cost and efficiency are
\begin{align}
\label{eq:transport-cost}
\mathrm{transportCost} &=
\sum_{\mathrm{edges}} \frac{Q_{ij}^2 L_{ij}}{D_{ij}+\epsilon},\\
\label{eq:efficiency}
\mathrm{efficiency} &=
\frac{\mathrm{deliveredNutrients}}
{\mathrm{transportCost}+\epsilon}.
\end{align}

\section{Numerical Methods}
The simulation uses fixed substeps inside each animation frame. Trail diffusion
uses explicit finite differences with stability substepping. Pressure is solved
approximately by Gauss-Seidel iterations with one pressure fixed to remove the
null space. Dijkstra is run on the current low-resolution graph. Randomness comes
from a deterministic xorshift generator seeded by the UI.

\section{Limitations and Future Work}
The model uses point particles and a scalar trail field rather than a deformable
continuous organism. Food scent is rebuilt from sources rather than diffused as a
separate PDE. The network graph is sampled from grid points, so ridge extraction
and branch-point detection are approximate. Future work could add WebGL rendering,
obstacles with richer geometry, persistent graph conductivity across rebuilds,
multi-species competition, better tube-width visualization, and calibrated
biophysical parameters.

\section{References}
\nocite{*}
\bibliographystyle{plain}
\bibliography{references}

\end{document}
